\textbf{Lösung zu Klausurvorbereitung 7}

\begin{loesung}
\textbf{1. Lösungsweg:} Satz von Euler, Schnelles Potenzieren
\begin{align*}
& e\,d \equiv 1 \pmod{\varphi(m)}, \quad \varphi(101) = 100, && a^{\varphi(m)} \equiv 1 \pmod m\\
\Rightarrow\quad & d = e^{\varphi(\varphi(m)) - 1} \equiv 1 \pmod{\varphi(m)} && \varphi(\varphi(m)) = \varphi(100) = 40\\
\Rightarrow\quad & \varphi(\varphi(m)) - 1 = 39 = 1 + 2 + 4 + 32\\
& 13,\quad 13^2 \equiv 69,\quad 13^4 \equiv 69^2 \equiv 61,\quad \dots,\quad 13^{32} \equiv 81\\
\Rightarrow\quad & d = e^{39} \equiv 1 \pmod{100}\\
\Rightarrow\quad & d \equiv 13 \cdot 69 \cdot 61 \cdot 81 \equiv 77 \pmod{100}
\end{align*}
% Hinweis: In der Quelle steht "13^2", "13^4", "69^2", "13^32" in Textschreibweise und "*" als Malzeichen.

\textbf{2. Lösungsweg:} Lemma von Bezout, EA
\[
e\,d - \varphi(m)\,y = 1 \quad\Leftrightarrow\quad e\,d \equiv 1 \pmod{\varphi(m)}
\]
EA mit $\varphi(m)$ und $e$
\begin{align*}
100 &= 7 \cdot 13 + 9\\
13 &= 1 \cdot 9 + 4\\
9 &= 2 \cdot 4 + 1\\
4 &= 4 \cdot 1
\end{align*}
\[
Q = \begin{pmatrix} 7 & 1\\ 1 & 0\end{pmatrix}
\begin{pmatrix} 1 & 1\\ 1 & 0\end{pmatrix}
\begin{pmatrix} 2 & 1\\ 1 & 0\end{pmatrix}
\begin{pmatrix} 4 & 1\\ 1 & 0\end{pmatrix}
= \begin{pmatrix} 8 & 7\\ 1 & 1\end{pmatrix}
\begin{pmatrix} 9 & 2\\ 4 & 1\end{pmatrix}
= \begin{pmatrix} 100 & 23\\ 13 & 3\end{pmatrix}
\]
$\operatorname{Det} Q = 100 \cdot 3 - 13 \cdot 23 = 1$\\
$\Rightarrow\quad d = -23 \equiv 77 \pmod{100}$
\end{loesung}
